\begin{blocksection}
\question Implement \texttt{reduce}, a tail recursive function that applies a
two-argument function \texttt{f} on the elements of a list from left to right,
starting with \texttt{start}.
\newline
\begin{lstlisting}
pub fn reduce(l: List(a), f: fn(a, a) -> a, start: a) -> a {
  case l {
    [] -> ______
    [first] -> ______________________________
    [first, ..rest] -> ______________________________
  }
}
\end{lstlisting}
\begin{solution}
\begin{lstlisting}
pub fn reduce(l: List(a), f: fn(a, a) -> a, start: a) -> a {
  case l {
    [] -> start
    [first] -> f(start, first)
    [first, ..rest] -> reduce(rest, f, f(start, first))
  }
}
\end{lstlisting}
\end{solution}

\begin{questionmeta}
  The purpose of this question is to demonstrate the use of generics in Gleam. \lstinline{reduce} accepts a list of any type, and only
  enforces that the type of the elements of the list is the same as the type of \lstinline{f} (both the arguments and return value), \lstinline{start},
  and the return type of \lstinline{reduce}. \lstinline{reduce} is tail recursive because the only case where a recursive call is made makes no
  modification to the return value of the recursive call.
\end{questionmeta}
\end{blocksection}